\documentclass[a4,11pt]{aleph-notas}

% -- Paquetes adicionales
\usepackage{enumitem}
\usepackage{aleph-comandos}
\usepackage{systeme}

% -- Datos
\institucion{Facultad de Ciencias Exactas, Naturales y Ambientale}
\carrera{Catálogo STEM}
\asignatura{Métodos Numéricos}
\tema{Ejercicios resueltos: Sistemas de ecuaciones no lineales}
\autor{Mario Cueva - Andrés Merino}
\fecha{Periodo 2026-1}

\logouno[0.18\textwidth]{Logos/logoPUCE_04_ac}
\definecolor{colortext}{HTML}{1957d1}
\definecolor{colordef}{HTML}{1957d1}
\fuente{montserrat}

% -- Comandos adicionales
\setlist[enumerate]{label=\roman*.}

\begin{document}

\encabezado

\section{Sistemas de ecuaciones no lineales}


%%%%%%%%%%%%%%%%%%%%%%%%%%%%%%%%%%%%%%%%
\begin{ejer}
    Dado el sistema de ecuaciones no lineales
    \[
        \begin{cases}
            f_1(x,y) = x^2 - 2x - y + \dfrac{1}{2} = 0,\\[6pt]
            f_2(x,y) = x^2 + 4y^2 - 4 = 0,
        \end{cases}
    \]
    \begin{enumerate}
        \item determinar la matriz jacobiana $J(x,y)$ del sistema;
        \item aplicar tres iteraciones del \textbf{método de Newton-Raphson} con la
              aproximación inicial $\mathbf{x}^{(0)}=\!\left(2,\;0{,}25\right)^{\!\top}$, y
        \item construir una tabla con la norma del residuo
              $\bigl\|F\!\left(\mathbf{x}^{(k)}\right)\bigr\|_\infty$ en cada iteración
              para evidenciar la convergencia cuadrática del método.
    \end{enumerate}
\end{ejer}

\begin{proof}[Solución]\hspace{0pt}

\textbf{i. Matriz jacobiana.}
Para $f_1(x,y)=x^2-2x-y+\tfrac{1}{2}$ y $f_2(x,y)=x^2+4y^2-4$:
\[
    J(x,y)
    = \begin{pmatrix}
        \dfrac{\partial f_1}{\partial x} & \dfrac{\partial f_1}{\partial y}\\[8pt]
        \dfrac{\partial f_2}{\partial x} & \dfrac{\partial f_2}{\partial y}
      \end{pmatrix}
    = \begin{pmatrix} 2x-2 & -1 \\ 2x & 8y \end{pmatrix}.
\]

\textbf{ii. Iteraciones de Newton-Raphson.}
En cada paso se resuelve el sistema lineal
\[
    J\!\left(\mathbf{x}^{(k)}\right)\boldsymbol{\delta}^{(k)}
    = -F\!\left(\mathbf{x}^{(k)}\right)
\]
y se actualiza $\mathbf{x}^{(k+1)}=\mathbf{x}^{(k)}+\boldsymbol{\delta}^{(k)}$.

\medskip
\textit{Iteración 1} ($k=0$, $\mathbf{x}^{(0)}=(2;\;0{,}25)^\top$):
\[
    F\!\left(\mathbf{x}^{(0)}\right)
    = \begin{pmatrix} 4-4-0{,}25+0{,}5\\ 4+4(0{,}0625)-4 \end{pmatrix}
    = \begin{pmatrix} 0{,}25\\ 0{,}25 \end{pmatrix},
    \qquad
    J\!\left(\mathbf{x}^{(0)}\right)
    = \begin{pmatrix} 2 & -1\\ 4 & 2 \end{pmatrix}.
\]
Resolviendo $J\boldsymbol{\delta}=-F$:
\[
    \begin{cases} 2\delta_1-\delta_2=-0{,}25,\\[3pt] 4\delta_1+2\delta_2=-0{,}25, \end{cases}
    \quad\implies\quad
    \delta_1=-0{,}09375,\quad
    \delta_2=\phantom{-}0{,}06250.
\]
\[
    \mathbf{x}^{(1)}
    = (2-0{,}09375;\;0{,}25+0{,}06250)^\top
    = (1{,}90625;\;0{,}31250)^\top.
\]

\textit{Iteración 2} ($k=1$, $\mathbf{x}^{(1)}=(1{,}90625;\;0{,}31250)^\top$):
\[
    F\!\left(\mathbf{x}^{(1)}\right)
    \approx \begin{pmatrix} 0{,}008789\\ 0{,}024414 \end{pmatrix},
    \qquad
    J\!\left(\mathbf{x}^{(1)}\right)
    = \begin{pmatrix} 1{,}81250 & -1\\ 3{,}81250 & 2{,}50000 \end{pmatrix}.
\]
Resolviendo $J\boldsymbol{\delta}=-F$:
\[
    \delta_1 \approx -0{,}005558,\qquad
    \delta_2 \approx -0{,}001284.
\]
\[
    \mathbf{x}^{(2)} \approx (1{,}900692;\;0{,}311216)^\top.
\]

\textit{Iteración 3} ($k=2$, $\mathbf{x}^{(2)}\approx(1{,}900692;\;0{,}311216)^\top$):
\[
    F\!\left(\mathbf{x}^{(2)}\right)
    \approx \begin{pmatrix} 3{,}0\times10^{-5}\\ 5{,}0\times10^{-5} \end{pmatrix}.
\]
La corrección $\boldsymbol{\delta}^{(2)}$ es del orden de $10^{-8}$;
el proceso ha convergido a $\mathbf{x}^*\approx(1{,}90068;\;0{,}31122)^\top$.

\medskip
\textbf{iii. Tabla de convergencia.}
\[
    \begin{array}{c|c|c}
        k & \mathbf{x}^{(k)} & \bigl\|F\!\left(\mathbf{x}^{(k)}\right)\bigr\|_\infty \\ \hline
        0 & (2{,}00000;\;0{,}25000)^\top & 2{,}50\times10^{-1} \\
        1 & (1{,}90625;\;0{,}31250)^\top & 2{,}44\times10^{-2} \\
        2 & (1{,}90069;\;0{,}31122)^\top & 5{,}0\phantom{0}\times10^{-5} \\
    \end{array}
\]
La caída de tres órdenes de magnitud entre las iteraciones 1 y 2 confirma la
\textbf{convergencia cuadrática} característica del método de Newton-Raphson.

\end{proof}


%%%%%%%%%%%%%%%%%%%%%%%%%%%%%%%%%%%%%%%%
\begin{ejer}
    Considerar el sistema de ecuaciones no lineales de tres variables:
    \[
        \begin{cases}
            f_1(x,y,z) = 3x - \cos(yz) - \dfrac{1}{2} = 0,\\[8pt]
            f_2(x,y,z) = x^2 - 81(y+0{,}1)^2 + \sin z + 1{,}06 = 0,\\[8pt]
            f_3(x,y,z) = e^{-xy} + 20z + \dfrac{10\pi-3}{3} = 0.
        \end{cases}
    \]
    \begin{enumerate}
        \item Reescribir el sistema en la forma de punto fijo $\mathbf{x}=G(\mathbf{x})$,
              despejando cada variable de su ecuación correspondiente.
        \item Aplicar cuatro iteraciones del \textbf{método de iteración de punto fijo}
              con la aproximación inicial
              $\mathbf{x}^{(0)}=(0{,}1;\;0{,}1;\;-0{,}1)^\top$.
        \item Verificar que la solución exacta del sistema es
              $\mathbf{x}^*=\!\left(\dfrac{1}{2},\;0,\;-\dfrac{\pi}{6}\right)^{\!\top}$.
    \end{enumerate}
\end{ejer}

\begin{proof}[Solución]\hspace{0pt}

\textbf{i. Forma de punto fijo.}
Despejando cada variable directamente:
\begin{align*}
    x &= g_1(x,y,z) = \dfrac{1}{3}\!\left(\cos(yz)+\dfrac{1}{2}\right),\\[6pt]
    y &= g_2(x,y,z) = \dfrac{1}{9}\sqrt{x^2+\sin z+1{,}06\,}-0{,}1,\\[6pt]
    z &= g_3(x,y,z) = \dfrac{1}{20}\!\left(-e^{-xy}-\dfrac{10\pi-3}{3}\right).
\end{align*}
Se elige la raíz positiva en $g_2$, adecuada para la región de interés $y\geq -0{,}1$.

\medskip
\textbf{ii. Iteraciones de punto fijo} ($\mathbf{x}^{(k+1)}=G\!\left(\mathbf{x}^{(k)}\right)$):
\[
    \begin{array}{c|ccc}
        k & x^{(k)} & y^{(k)} & z^{(k)} \\ \hline
        0 & 0{,}10000 & 0{,}10000 & -0{,}10000 \\
        1 & 0{,}49999 & 0{,}00944 & -0{,}52310 \\
        2 & 0{,}50000 & 0{,}00001 & -0{,}52336 \\
        3 & 0{,}50000 & 0{,}00000 & -0{,}52360 \\
        4 & 0{,}50000 & 0{,}00000 & -0{,}52360 \\
    \end{array}
\]
La sucesión converge rápidamente a
$\mathbf{x}^*=(0{,}5;\;0;\;-\pi/6)^\top$, con $-\pi/6\approx-0{,}52360$.

\medskip
\textbf{iii. Verificación de la solución exacta.}
Sustituyendo $\mathbf{x}^*=\!\left(\tfrac{1}{2},\,0,\,-\tfrac{\pi}{6}\right)^\top$:
\begin{align*}
    f_1 &= 3\cdot\frac{1}{2}-\cos\!\left(0\cdot\!\left(-\frac{\pi}{6}\right)\right)-\frac{1}{2}
         = \frac{3}{2}-\cos(0)-\frac{1}{2}
         = \frac{3}{2}-1-\frac{1}{2}=0.\\[8pt]
    f_2 &= \left(\frac{1}{2}\right)^{\!2}-81(0{,}1)^2+\sin\!\left(-\frac{\pi}{6}\right)+1{,}06
         = 0{,}25-0{,}81-0{,}50+1{,}06=0.\\[8pt]
    f_3 &= e^{-\frac{1}{2}\cdot 0}+20\!\left(-\frac{\pi}{6}\right)+\frac{10\pi-3}{3}
         = 1-\frac{10\pi}{3}+\frac{10\pi}{3}-1=0.
\end{align*}
Queda verificado que $\mathbf{x}^*$ es solución exacta del sistema.

\end{proof}


\end{document}
